\secnumbersection{TRABAJO RELACIONADO}

El estudio de la movilidad humana y la predicción de flujos origen-destino (OD) ha experimentado una evolución metodológica fundamental en las últimas décadas. Este capítulo no busca enumerar exhaustivamente todas las investigaciones previas, sino trazar una línea evolutiva que contraste las fortalezas y debilidades de los distintos enfoques adoptados en la literatura científica. A través de este análisis, se fundamentarán las decisiones metodológicas que sustentan la presente memoria, demostrando por qué la integración de modelos de aprendizaje profundo con datos geoespaciales abiertos es la alternativa más adecuada para el contexto actual.

% ----------------------------------------------------------------
\subsection{Descubrimiento de patrones en la movilidad individual}
\label{sec:patrones_movilidad}
% ----------------------------------------------------------------

Durante mucho tiempo, la falta de datos de trayectoria masivos obligó a los investigadores a extrapolar los patrones de movilidad humana a partir de la dispersión de billetes de banco, asumiendo que los desplazamientos individuales seguían trayectorias puramente aleatorias y no acotadas, similares a los vuelos de Lévy (\textit{Lévy flights}). 

Este paradigma fue refutado por Gonz\'alez et al. \cite{gonzalez2008understanding}, quienes utilizaron registros masivos de telefonía móvil (CDR, por sus siglas en inglés) de 100.000 usuarios durante seis meses para analizar la movilidad humana a nivel microscópico. Su trabajo demostró que, a diferencia del movimiento browniano o de los vuelos de Lévy, las trayectorias humanas son altamente predecibles, regulares y están fuertemente confinadas espacialmente. Los autores introdujeron el concepto de \textit{radio de giro} para cuantificar la distancia típica recorrida por un individuo, demostrando que la aparente distribución pesada de saltos poblacionales es, en realidad, el resultado de combinar trayectorias individuales acotadas con una gran heterogeneidad en los radios de giro de las personas.

Esta contribución fue clave porque estableció que la movilidad urbana no es un fenómeno caótico, sino un sistema gobernado por rutinas (hogar-trabajo) y patrones espaciales predecibles. Esta previsibilidad estadística es la premisa fundamental que justifica el desarrollo de modelos matemáticos y algorítmicos para la predicción de flujos agregados.

% ----------------------------------------------------------------
\subsection{Unificación de modelos clásicos y enfoques continuos}
\label{sec:modelos_clasicos_continuos}
% ----------------------------------------------------------------

Con la confirmación de la regularidad en la movilidad humana, los esfuerzos se centraron en predecir los volúmenes de viaje entre distintas zonas de una ciudad o país. Tradicionalmente, este problema se abordó mediante el modelo gravitatorio \cite{zipf1946p} y el modelo de oportunidades intervinientes \cite{stouffer1940intervening}. Más tarde, el modelo de radiación \cite{simini2012universal} ofreció una solución libre de parámetros que superaba varias limitaciones empíricas de la formulación gravitatoria.

Sin embargo, estos modelos presentaban un problema matemático estructural: dependían fuertemente de la partición espacial utilizada (comunas, distritos, censos), lo que generaba inconsistencias espaciales. Además, salvo casos triviales, las formulaciones discretas como el modelo gravitatorio violaban la propiedad de \textit{aditividad} (el flujo hacia una región compuesta no siempre equivalía a la suma de los flujos hacia sus subregiones).

Para resolver esto, Simini, Maritan y N\'eda \cite{simini2013human} propusieron un marco unificado en el continuo. Los autores formularon un principio estocástico general de selección de beneficios laborales que no solo preservaba estrictamente la aditividad de los flujos, sino que demostró que la gravedad, las oportunidades intervinientes y la radiación eran, matemáticamente, casos particulares de una misma familia de procesos de selección espacial continua. Además, introdujeron un parámetro de selección ($\lambda$) que flexibilizó el modelo de radiación para adaptarlo a diferentes escalas geográficas.

A pesar de su solidez teórica y elegancia matemática, este paradigma continuo arrastra una limitación fundamental de representación: asume que la cantidad de oportunidades en una región es estrictamente proporcional a su masa poblacional. En la realidad urbana, existe una profunda segregación funcional del suelo (zonas industriales, distritos financieros, áreas puramente residenciales) que rompe esta proporcionalidad lineal, exigiendo modelos capaces de entender el tejido urbano más allá de la simple aglomeración de personas.

% ----------------------------------------------------------------
\subsection{Deep Learning y modelos espacio-temporales en movilidad}
\label{sec:deep_learning_movilidad}
% ----------------------------------------------------------------

Para superar las limitaciones lineales y la dependencia exclusiva de la densidad poblacional, la literatura reciente ha migrado hacia el aprendizaje automático (\textit{Machine Learning}) y el aprendizaje profundo (\textit{Deep Learning}) \cite{liu2024interdisciplinary}. Las arquitecturas neuronales han demostrado una capacidad superior para capturar la no linealidad del entorno urbano.

Un avance significativo en esta línea es la incorporación explícita de Puntos de Interés (POIs) y características del entorno construido como variables condicionantes. Por ejemplo, Luca et al. \cite{luca2024tsmob} propusieron el \textit{framework} TS-Mob, que combina modelos fundacionales de series de tiempo (\textit{TimesFM}) con un módulo de atractividad gravitacional basado en perceptrones multicapa (MLP). Este modelo integra 16 categorías de POIs extraídas de bases de datos abiertas, permitiendo modular la probabilidad de atracción estructural de un destino en función de la verdadera oferta de servicios comerciales, educativos o de salud, y no solo de su población.

No obstante, la transición hacia arquitecturas profundas introdujo nuevos desafíos. Las redes neuronales profundas y espaciales requieren enormes volúmenes de datos para evitar el sobreajuste y, más críticamente, operan como "cajas negras" \cite{liu2024interdisciplinary}. Aunque reducen drásticamente el error absoluto medio (MAE) y mejoran la similitud estructural de flujos (CPC) respecto a los modelos clásicos, la incapacidad de explicar por qué el modelo predice un flujo determinado reduce su adopción institucional, un aspecto crítico en la planificación urbana pública.

% ----------------------------------------------------------------
\subsection{Generación de flujos mediante datos no convencionales}
\label{sec:flujos_no_convencionales}
% ----------------------------------------------------------------

En el estado del arte más reciente (2024-2025), la investigación ha explorado métodos para mitigar la dependencia de datos socioeconómicos hiperlocales o de costosas encuestas de transporte. Destaca el paradigma de inferencia mediante imágenes satelitales y modelos generativos de grafos.

Rong et al. \cite{rong2024satellites} desarrollaron \textit{GlODGen}, un generador global de flujos OD basado en modelos de difusión (\textit{Denoising Diffusion Networks}). Este \textit{framework} utiliza un codificador visual multimodal (\textit{RemoteCLIP}) para extraer un \textit{embedding} semántico de 1.024 dimensiones directamente desde imágenes satelitales de alta resolución de la ciudad. Los autores demostraron que el aspecto visual del entorno urbano (visto desde el espacio) y un mapa global de población poseen una expresividad superior al 98\% en comparación con el uso de datos multisectoriales complejos (catastros, POIs, atributos socioeconómicos). Su modelo, entrenado en Estados Unidos, logró transferibilidad a ciudades de Europa, China, Brasil y África.

A pesar del salto tecnológico que representan estos modelos generativos satelitales, presentan ciertas barreras prácticas. Suelen estar condicionados por sesgos geográficos (morfología urbana de los países donde fueron entrenados) y su desempeño es sensible a la calidad, iluminación y nubosidad de las imágenes satelitales. Además, los \textit{embeddings} visuales de 1.024 dimensiones sacrifican completamente la interpretabilidad semántica: es imposible aislar el impacto de variables específicas como "presencia de escuelas" o "red de hospitales" sobre la predicción del flujo, variables que son fundamentales para simular escenarios de impacto urbano.

% ----------------------------------------------------------------
\subsection{En Resumen}
% ----------------------------------------------------------------

La revisión de la literatura evidencia un claro \textit{trade-off} en la evolución de los modelos de movilidad:

\begin{itemize}
    \item Los modelos físicos y estadísticos clásicos (Secciones \ref{sec:patrones_movilidad} y \ref{sec:modelos_clasicos_continuos}) ofrecen transparencia matemática, rigor explicativo y garantías de aditividad, pero fracasan al intentar modelar las complejidades y no linealidades del uso de suelo moderno al depender exclusivamente de la distancia y la masa poblacional como variables \textit{proxy} de oportunidad.
    \item Los modelos avanzados basados en Deep Learning, series de tiempo temporales e imágenes satelitales (Secciones \ref{sec:deep_learning_movilidad} y \ref{sec:flujos_no_convencionales}) alcanzan una alta precisión predictiva y transferibilidad global, pero se estructuran como cajas negras inflexibles, opacas y computacionalmente costosas que limitan su utilidad para fundamentar políticas públicas explicables.
\end{itemize}

En este contexto metodológico, el uso del modelo \textit{Deep Gravity} \cite{simini2021deep} alimentado con datos de OpenStreetMap (OSM) y analizado mediante técnicas de explicabilidad SHAP \cite{lundberg2017shap} se posiciona como una solución de equilibrio o \textit{sweet spot} para el caso de Santiago de Chile. 

Por un lado, \textit{Deep Gravity} reemplaza la rigidez heurística de los exponentes gravitacionales por una red neuronal capaz de aprender no linealidades (superando a los clásicos), y lo hace utilizando variables explícitas y semánticamente claras de la estructura espacial provistas por OSM (centros de salud, red vial, escuelas, comercio), superando la ambigüedad de los \textit{embeddings} satelitales. Por otro lado, al preservar una estructura arquitectónica relativamente simple, facilita la aplicación eficiente de métodos de Explicabilidad de la IA (como SHAP), lo que permitirá no solo estimar los flujos de la ciudad con alta precisión, sino auditar \textit{por qué} ocurren, proveyendo evidencia interpretable para la planificación territorial en el Gran Santiago.
